\exercise{8.6}{Absolute values at zero and an alternating counterexample}

\goal Prove (a) $s_n\to0$ if and only if $\abs{s_n}\to0$, and
(b) the sequence $s_n=(-1)^n$ has convergent absolute values but is
itself divergent.

Throughout these exercises, convergence is used in the precise form of
Definition~7.1: $x_n\to L$ means that for every $\varepsilon>0$
there is a real number $N$ such that
\[
 n\in\N,\quad n>N\quad\Longrightarrow\quad
 \abs{x_n-L}<\varepsilon.
\]
A threshold $N$ need not be an integer. Whenever an index above a
specified threshold is needed, the Archimedean property supplies a
positive integer above it.

\step{Part (a), Step 1: compare the two errors exactly.}
For every $n\in\N$, the number $\abs{s_n}$ is nonnegative. Therefore
\[
 \abs{\abs{s_n}-0}=\abs{\abs{s_n}}=\abs{s_n}
                =\abs{s_n-0}.
\]
The errors in the two proposed zero limits are the same real number,
not merely bounded by each other.

\step{Part (a), Step 2: prove the forward implication.}
Assume $s_n\to0$, and let $\varepsilon>0$ be arbitrary.
The definition supplies $N$ such that
\[
 n>N\quad\Longrightarrow\quad\abs{s_n-0}<\varepsilon.
\]
For each such index, the identity in Step~1 gives
\[
 \abs{\abs{s_n}-0}=\abs{s_n-0}<\varepsilon.
\]
Thus the same threshold $N$ verifies convergence of $\abs{s_n}$ to $0$.

\step{Part (a), Step 3: prove the reverse implication.}
Assume $\abs{s_n}\to0$, and let $\varepsilon>0$ be arbitrary.
Choose $N$ from this convergence so that
\[
 n>N\quad\Longrightarrow\quad\abs{\abs{s_n}-0}<\varepsilon.
\]
Using Step~1 again yields
\[
 \abs{s_n-0}=\abs{\abs{s_n}-0}<\varepsilon
 \qquad(n>N).
\]
This is the definition of $s_n\to0$.

\answer{Both implications have been proved with the same error and
the same threshold:
\[
 \boxed{s_n\to0\quad\Longleftrightarrow\quad\abs{s_n}\to0.}
\]}

\step{Part (b), Step 1: prove convergence of the absolute values.}
If $s_n=(-1)^n$, then $s_n$ equals either $1$ or $-1$. Hence
$\abs{s_n}=1$ for every $n$.
For any $\varepsilon>0$, choose $N=0$. Every positive integer $n>N$
satisfies
\[
 \abs{\abs{s_n}-1}=\abs{1-1}=0<\varepsilon.
\]
Therefore $\abs{s_n}\to1$.

\step{Part (b), Step 2: rule out every possible limit of the original sequence.}
Suppose, for a contradiction, that $s_n\to L$ for some real $L$.
Apply the definition with the fixed positive tolerance $\varepsilon=1$.
There is a real $N$ such that
\[
 n>N\quad\Longrightarrow\quad\abs{s_n-L}<1.
\]
Choose an integer $k>\max\{N,1\}$. Then both $2k$ and $2k+1$ are
positive integers larger than $N$. Their terms are
\[
 s_{2k}=1,\qquad s_{2k+1}=-1.
\]
The triangle inequality gives
\begin{align*}
 2
   &=\abs{1-(-1)}\\
   &=\abs{(1-L)+(L-(-1))}\\
   &\le\abs{1-L}+\abs{L-(-1)}\\
   &=\abs{s_{2k}-L}+\abs{s_{2k+1}-L}\\
   &<1+1=2.
\end{align*}
The equality of the absolute values in the fourth line uses
$\abs{-x}=\abs{x}$. The resulting strict inequality $2<2$ is impossible.

\answer{The absolute-value sequence converges to $1$, but
$((-1)^n)$ has no real limit. Thus existence of a limit for
$\abs{s_n}$ alone does not guarantee existence of a limit for $s_n$.}

\exercise{8.7}{Three divergent sequences}

\goal Show that none of the following sequences converges:
\[
 \cos\left(\frac{n\pi}{3}\right),\qquad
 (-1)^n n,\qquad
 \sin\left(\frac{n\pi}{3}\right).
\]

\idea We will use only the definition of convergence and the triangle
inequality. If a sequence converged to $L$, then for any chosen
$\varepsilon>0$, all sufficiently late terms would be within
$\varepsilon$ of $L$. Any two such terms would satisfy
\[
 \abs{s_p-s_q}
  =\abs{(s_p-L)+(L-s_q)}
  \le\abs{s_p-L}+\abs{s_q-L}
  <2\varepsilon.
\]
For each proposed sequence, we will produce two arbitrarily late terms
whose distance contradicts this inequality. This displayed observation
has just been proved directly, so no theorem about subsequences or
bounded convergent sequences is needed.

\step{Part (a), Step 1: select a tolerance under an assumed limit.}
Suppose
\[
 s_n=\cos\left(\frac{n\pi}{3}\right)\to L.
\]
With $\varepsilon=1$, choose $N$ such that every positive integer
$n>N$ satisfies $\abs{s_n-L}<1$.

\step{Part (a), Step 2: find late terms equal to $1$ and $-1$.}
Choose an integer $k>\max\{N,1\}$, and put
\[
 p=6k,\qquad q=6k+3.
\]
Both indices exceed $N$. The period $2\pi$ of cosine gives
\[
 s_p=\cos(2k\pi)=1,\qquad
 s_q=\cos(2k\pi+\pi)=-1.
\]
Their distance is therefore $\abs{s_p-s_q}=2$.
The two-late-terms estimate gives
\[
 2=\abs{s_p-s_q}
   \le\abs{s_p-L}+\abs{s_q-L}
   <1+1=2,
\]
a contradiction.

\answer{The sequence $\bigl(\cos(n\pi/3)\bigr)$ does not converge.}

\step{Part (b), Step 1: use the definition with a fixed tolerance.}
Suppose $s_n=(-1)^n n\to L$.
Choose $N$ from the definition with $\varepsilon=1$.

\step{Part (b), Step 2: compute the distance between late even and odd terms.}
Choose an integer $k>\max\{N,1\}$.
Then $p=2k$ and $q=2k+1$ are larger than $N$ and
\[
 s_p=2k,\qquad s_q=-(2k+1).
\]
Their difference is
\[
 s_p-s_q=2k+(2k+1)=4k+1>2.
\]
Since this difference is positive, its absolute value is also $4k+1$.
But convergence would give
\[
 4k+1
 =\abs{s_p-s_q}
 \le\abs{s_p-L}+\abs{s_q-L}
 <1+1=2.
\]
This contradicts $4k+1>2$.

\answer{The sequence $\bigl((-1)^n n\bigr)$ does not converge.
The proof uses two explicit terms in every sufficiently late tail,
rather than assuming a boundedness theorem from a later section.}

\step{Part (c), Step 1: choose a tolerance suited to the sine values.}
Suppose
\[
 s_n=\sin\left(\frac{n\pi}{3}\right)\to L.
\]
Choose the positive tolerance $\varepsilon=\sqrt3/2$.
The definition gives $N$ such that $n>N$ implies
$\abs{s_n-L}<\sqrt3/2$.

\step{Part (c), Step 2: use two positions in the six-term period.}
Choose an integer $k>\max\{N,1\}$ and let
\[
 p=6k+1,\qquad q=6k+4.
\]
Both are larger than $N$. By the period $2\pi$ of sine,
\[
 s_p=\sin\left(2k\pi+\frac{\pi}{3}\right)=\frac{\sqrt3}{2},
 \qquad
 s_q=\sin\left(2k\pi+\frac{4\pi}{3}\right)=-\frac{\sqrt3}{2}.
\]
In particular, $\abs{s_p-s_q}=\sqrt3$.
The convergence estimates would imply
\[
 \sqrt3
 =\abs{s_p-s_q}
 \le\abs{s_p-L}+\abs{s_q-L}
 <\frac{\sqrt3}{2}+\frac{\sqrt3}{2}
 =\sqrt3.
\]
This is impossible.

\answer{The sequence $\bigl(\sin(n\pi/3)\bigr)$ does not converge.
Every sequence in the exercise has now been ruled out by the
definition of convergence itself.}

\exercise{8.8}{Three radical limits with explicit error bounds}

\goal Prove the three stated limits:
\begin{align*}
 \sqrt{n^2+1}-n&\longrightarrow0,\\
 \sqrt{n^2+n}-n&\longrightarrow\frac12,\\
 \sqrt{4n^2+n}-2n&\longrightarrow\frac14.
\end{align*}
The index $n$ is always a positive integer. In particular, division by
$n$ is legitimate, and $\sqrt{n^2}=n$ rather than $-n$.

\idea Multiplication by a conjugate removes the subtraction of large
radicals. We will then bound each \emph{error from its claimed limit}
by a positive multiple of $1/n$. The estimates below are ordinary
algebraic identities and inequalities, not applications of limit laws.

\step{Part (a), Step 1: rationalize the expression and check its sign.}
Let $u_n=\sqrt{n^2+1}-n$.
Because $n\ge1$, both $n$ and $\sqrt{n^2+1}$ are positive.
Moreover, $\sqrt{n^2+1}>n$: otherwise squaring a supposed inequality
$\sqrt{n^2+1}\le n$ between nonnegative numbers would give
$n^2+1\le n^2$, a contradiction. Thus $u_n>0$.
Its conjugate $\sqrt{n^2+1}+n$ is positive, so
\begin{align*}
 u_n
  &=\frac{(\sqrt{n^2+1}-n)(\sqrt{n^2+1}+n)}
          {\sqrt{n^2+1}+n}\\
  &=\frac{(n^2+1)-n^2}{\sqrt{n^2+1}+n}\\
  &=\frac1{\sqrt{n^2+1}+n}.
\end{align*}

\step{Part (a), Step 2: obtain an error estimate valid for every index.}
The denominator satisfies $\sqrt{n^2+1}+n\ge2n>0$.
Taking reciprocals of positive quantities reverses the comparison, so
\[
 \abs{u_n-0}=u_n\le\frac1{2n}.
\]

\step{Part (a), Step 3: choose and verify a threshold.}
Let $\varepsilon>0$ be arbitrary, and choose
\[
 N=\frac1{2\varepsilon}.
\]
For every positive integer $n>N$, multiplying by $2\varepsilon>0$
gives $2\varepsilon n>1$. Division by $2n>0$ gives
$1/(2n)<\varepsilon$. Hence
\[
 n>N\quad\Longrightarrow\quad
 \abs{u_n-0}\le\frac1{2n}<\varepsilon.
\]
This is precisely $u_n\to0$.

\answer{\(\boxed{\displaystyle\lim_{n\to\infty}
(\sqrt{n^2+1}-n)=0}\). An explicit error bound is $1/(2n)$.}

\step{Part (b), Step 1: rationalize and factor out the positive index.}
Let $v_n=\sqrt{n^2+n}-n$. Multiplying by its positive conjugate gives
\[
 v_n=\frac{(n^2+n)-n^2}{\sqrt{n^2+n}+n}
     =\frac{n}{\sqrt{n^2+n}+n}.
\]
Set
\[
 r_n=\sqrt{1+\frac1n}.
\]
The quantity $nr_n$ is positive, and its square is
\[
 (nr_n)^2=n^2\left(1+\frac1n\right)=n^2+n.
\]
By uniqueness of the nonnegative square root,
$\sqrt{n^2+n}=nr_n$. Substitution gives the exact identity
\[
 v_n=\frac{n}{nr_n+n}
     =\frac1{r_n+1}.
\]
The number $r_n$ is greater than $1$, because $r_n^2=1+1/n>1$.
Thus $r_n+1>2$, and $0<v_n<1/2$.

\step{Part (b), Step 2: rationalize the error from $1/2$.}
The sign information in Step~1 gives
\begin{align*}
 \abs{v_n-\tfrac12}
  &=\frac12-\frac1{r_n+1}\\
  &=\frac{(r_n+1)-2}{2(r_n+1)}\\
  &=\frac{r_n-1}{2(r_n+1)}.
\end{align*}
A second conjugate calculation yields
\[
 r_n-1
   =\frac{(r_n-1)(r_n+1)}{r_n+1}
   =\frac{r_n^2-1}{r_n+1}
   =\frac{1/n}{r_n+1}.
\]
Therefore
\[
 \abs{v_n-\tfrac12}
   =\frac1{2n(r_n+1)^2}.
\]
Since $r_n\ge1$, we have $(r_n+1)^2\ge4$, so
\[
 \abs{v_n-\tfrac12}\le\frac1{8n}.
\]

\step{Part (b), Step 3: verify the definition of convergence.}
Let $\varepsilon>0$ be arbitrary and choose $N=1/(8\varepsilon)$.
For every positive integer $n>N$, multiplication and division by
positive quantities give $1/(8n)<\varepsilon$.
Consequently
\[
 \abs{v_n-\tfrac12}\le\frac1{8n}<\varepsilon
 \qquad(n>N).
\]
The definition proves the desired limit.

\answer{\(\boxed{\displaystyle\lim_{n\to\infty}
(\sqrt{n^2+n}-n)=\tfrac12}\). An explicit error bound is $1/(8n)$.}

\step{Part (c), Step 1: rationalize and identify a simpler denominator.}
Let $w_n=\sqrt{4n^2+n}-2n$.
Its conjugate is positive, and multiplication gives
\[
 w_n=\frac{(4n^2+n)-4n^2}{\sqrt{4n^2+n}+2n}
     =\frac{n}{\sqrt{4n^2+n}+2n}.
\]
Define $R_n=\sqrt{4+1/n}$. Since $n>0$, the quantity $nR_n$ is
the nonnegative square root of $4n^2+n$:
\[
 (nR_n)^2=n^2\left(4+\frac1n\right)=4n^2+n.
\]
Hence
\[
 w_n=\frac{n}{nR_n+2n}=\frac1{R_n+2}.
\]
Because $R_n^2=4+1/n>4$ and $R_n\ge0$, we have $R_n>2$.
Thus $R_n+2>4$ and $0<w_n<1/4$.

\step{Part (c), Step 2: compute and bound the error from $1/4$.}
Using the sign of the error, put the fractions over one denominator:
\begin{align*}
 \abs{w_n-\tfrac14}
   &=\frac14-\frac1{R_n+2}\\
   &=\frac{(R_n+2)-4}{4(R_n+2)}\\
   &=\frac{R_n-2}{4(R_n+2)}.
\end{align*}
The second rationalization gives
\[
 R_n-2
   =\frac{R_n^2-4}{R_n+2}
   =\frac{1/n}{R_n+2}.
\]
It follows that
\[
 \abs{w_n-\tfrac14}
   =\frac1{4n(R_n+2)^2}
   \le\frac1{64n},
\]
because $R_n\ge2$ implies $(R_n+2)^2\ge16$.

\step{Part (c), Step 3: give the final tolerance-dependent threshold.}
Let $\varepsilon>0$ be arbitrary and choose $N=1/(64\varepsilon)$.
Then $n>N$ implies $1/(64n)<\varepsilon$, so
\[
 \abs{w_n-\tfrac14}\le\frac1{64n}<\varepsilon
 \qquad(n\in\N,\ n>N).
\]
This verifies convergence directly.

\answer{\(\boxed{\displaystyle\lim_{n\to\infty}
(\sqrt{4n^2+n}-2n)=\tfrac14}\).
All three limits were proved by explicit inequalities, without
using algebraic limit theorems from Section~9.}

\exercise{8.9}{Eventual bounds pass to a convergent limit}

\goal Let $s_n\to s$. Prove:
\begin{enumerate}[label=(\alph*)]
 \item If $s_n\ge a$ for all but finitely many indices, then $s\ge a$.
 \item If $s_n\le b$ for all but finitely many indices, then $s\le b$.
 \item If $s_n\in[a,b]$ for all but finitely many indices, then
       $s\in[a,b]$.
\end{enumerate}

\step{A preliminary detail: replace finitely many exceptions by one threshold.}
If a property fails only on a finite set $E\subseteq\N$ of indices,
define $K=0$ if $E$ is empty and $K=\max E$ otherwise.
Every positive integer $n>K$ is outside $E$, so every such index
satisfies the property. This explains each use of an eventual-bound
threshold below.

\step{Part (a), Step 1: suppose the limit violates the eventual lower bound.}
Choose $K$ so that $n>K$ implies $s_n\ge a$.
Suppose, for a contradiction, that $s<a$. Then the gap $a-s$ is
positive. Set
\[
 \varepsilon=\frac{a-s}{2}>0.
\]
Convergence supplies $N_0$ such that $n>N_0$ implies
$\abs{s_n-s}<\varepsilon$.

\step{Part (a), Step 2: combine the two tail conditions.}
Choose a positive integer $n>\max\{K,N_0\}$.
The condition $n>K$ gives $s_n\ge a$. The condition $n>N_0$ gives
\[
 s_n-s<\varepsilon,\qquad
 s_n<s+\varepsilon=\frac{s+a}{2}.
\]
But
\[
 a-\frac{s+a}{2}=\frac{a-s}{2}>0.
\]
Thus $s_n<(s+a)/2<a$, contradicting $s_n\ge a$.

\answer{The inequality $s<a$ is impossible. Therefore
\(\boxed{s\ge a}\). Equality $s=a$ is permitted; for example, the
constant sequence $s_n=a$ has limit $a$ and satisfies every lower bound
in the hypothesis.}

\step{Part (b), Step 1: suppose the limit violates the eventual upper bound.}
Choose $K$ so that $n>K$ implies $s_n\le b$.
Suppose $s>b$, and choose the positive tolerance
\[
 \varepsilon=\frac{s-b}{2}.
\]
Convergence gives $N_0$ such that $n>N_0$ implies
$\abs{s_n-s}<\varepsilon$.

\step{Part (b), Step 2: choose one index after both thresholds.}
Take a positive integer $n>\max\{K,N_0\}$.
Then $s_n\le b$, whereas the convergence inequality gives
\[
 -\varepsilon<s_n-s,\qquad
 s_n>s-\varepsilon=\frac{s+b}{2}.
\]
The midpoint exceeds $b$ because
\[
 \frac{s+b}{2}-b=\frac{s-b}{2}>0.
\]
Hence $s_n>b$, contradicting $s_n\le b$.

\answer{The inequality $s>b$ is impossible. Therefore
\(\boxed{s\le b}\). The endpoint $s=b$ is allowed.}

\step{Part (c), Step 1: unpack membership of a closed interval.}
For $a\le b$,
\[
 x\in[a,b]\quad\Longleftrightarrow\quad a\le x\le b.
\]
If all but finitely many $s_n$ belong to this interval, then all but
finitely many satisfy $s_n\ge a$, and all but finitely many satisfy
$s_n\le b$. The finite-exception threshold may be chosen from the
indices where interval membership fails, so both inequalities hold
beyond the same threshold.

\step{Part (c), Step 2: apply both proved statements to the same limit.}
Part (a) gives $s\ge a$, and part (b) gives $s\le b$.
Together these say $a\le s\le b$, which is exactly $s\in[a,b]$.
If $a=b$, the two inequalities force $s=a=b$, so the singleton case
is included as well.

\answer{For a convergent sequence,
\[
 \boxed{s_n\in[a,b]\text{ for all but finitely many }n
          \quad\Longrightarrow\quad \lim s_n\in[a,b].}
\]
The finitely many exceptional early terms do not affect any of these
conclusions because an index can be chosen after every exception
and after the convergence threshold.}

\exercise{8.10}{A strict bound for the limit becomes an eventual strict bound}

\goal If $s_n\to s$ and $s>a$, prove that there is a real $N$ such that
\[
 n\in\N,\quad n>N\quad\Longrightarrow\quad s_n>a.
\]

\idea The positive gap $s-a$ allows us to place a whole
$\varepsilon$-neighborhood of the limit above $a$.

\step{Step 1: choose a tolerance smaller than the gap.}
Because $s>a$, the real number $s-a$ is positive. Define
\[
 \varepsilon=\frac{s-a}{2}>0.
\]
The tolerance is fixed from the limit and the bound; it does not depend
on the index $n$.

\step{Step 2: obtain a threshold from convergence.}
The definition of $s_n\to s$ gives a real number $N$ such that
\[
 n\in\N,\quad n>N\quad\Longrightarrow\quad
 \abs{s_n-s}<\varepsilon.
\]
For any index in this tail, the absolute-value inequality implies
\[
 -\varepsilon<s_n-s<\varepsilon.
\]
The left inequality gives $s_n>s-\varepsilon$.

\step{Step 3: check that the lower edge is above the desired bound.}
Substitution and simplification give
\[
 s-\varepsilon
   =s-\frac{s-a}{2}
   =\frac{2s-s+a}{2}
   =\frac{s+a}{2}.
\]
Also
\[
 \frac{s+a}{2}-a=\frac{s-a}{2}>0,
\]
so $(s+a)/2>a$. Combining these strict comparisons yields
\[
 s_n>s-\varepsilon=\frac{s+a}{2}>a
 \qquad(n\in\N,\ n>N).
\]

\answer{The threshold supplied for
$\varepsilon=(s-a)/2$ satisfies the required implication:
\[
 \boxed{n>N\quad\Longrightarrow\quad s_n>a.}
\]
The strict hypothesis $s>a$ matters: if $s=a$ and $s_n=a$ for all $n$,
the sequence converges to $a$ but none of its terms is greater than $a$.}
